\section*{Chapter \ref{ch:b1:constrained-optimisation} --- Optimisation under Constraint}
% Numbers recomputed in code/bachelor-1/02-constrained-optimisation/exercises/solutions.py (tested,
% with a numerical optimiser as a check).

\begin{solution}{exo:b1:constrained-optimisation:read-chart}
\begin{enumerate}
  \item Pay differs by a factor of $49.48/2.20 = 22.5$. Hours differ by
    $2064 - 1504 = 560$ a year, which is $560/2064 = 27\%$ fewer in the top
    tenth.
  \item $560/48 = 11.7$ hours a week fewer.
  \item No: each point is a different country, with its own taxes, its own
    mix of jobs and its own habits, not one worker facing two wages.
\end{enumerate}
\end{solution}

\begin{solution}{exo:b1:constrained-optimisation:substitution}
$y = (24 - 2x)/3$, so $xy = (24x - 2x^2)/3$, with derivative
$(24 - 4x)/3 = 0$ at $x = 6$; the second derivative $-4/3$ is negative.
Then $y = 4$ and $xy = 24$. Tangency: $f_x/f_y = y/x = 4/6 = 2/3$ and
$g_x/g_y = 2/3$.
\end{solution}

\begin{solution}{exo:b1:constrained-optimisation:level-slope}
$f_x = 2xy = 4$ and $f_y = x^2 = 1$ at $(1, 2)$, so the slope is $-4/1 = -4$.
The level is $f(1, 2) = 2$, and the curve is $x^2y = 2$, that is $y = 2/x^2$
(whose derivative $-4/x^3$ is indeed $-4$ at $x = 1$).
\end{solution}

\begin{solution}{exo:b1:constrained-optimisation:binding}
At $(6, 3)$, $x + 2y = 12$ holds with equality: it binds. $x = 6 < 8$ and
$y = 3 < 5$: both are slack, with \omterm{def:b1:constrained-optimisation:shadow-price}{shadow prices} of zero. With $x \le 4$, the
point $(6, 3)$ is no longer allowed; the best point keeps $x$ as close as
allowed to the old tangency, $x = 4$, and spends the rest on $y$:
$y = (12 - 4)/2 = 4 \le 5$. The solution is $(4, 4)$, with $f = 4$, and now two
constraints bind.
\end{solution}

\begin{solution}{exo:b1:constrained-optimisation:lagrange-log}
$\mathcal L = \ln x + 2\ln y - \lambda(p_1x + p_2y - m)$. The conditions are
$1/x = \lambda p_1$ and $2/y = \lambda p_2$. Dividing, $p_2 y = 2p_1 x$; the
constraint gives $3p_1x = m$. So
\[
x^* = \frac{m}{3p_1}, \qquad y^* = \frac{2m}{3p_2}, \qquad \lambda^* = \frac{1}{p_1 x^*} = \frac{3}{m}.
\]
At $p_1 = 1$, $p_2 = 2$, $m = 30$: $x^* = 10$, $y^* = 10$, $\lambda^* = 0.1$.
\end{solution}

\begin{solution}{exo:b1:constrained-optimisation:lambda-units}
\begin{enumerate}
  \item Units of the objective per dalen: $\lambda^*$ is the increase in the
    maximum value per extra dalen of $m$, the marginal value of income.
  \item $\Delta V \approx \lambda^* \Delta m = 0.1 \times 10 = 1.0$.
  \item $V(m) = \ln(m/3) + 2\ln(m/3) = 3\ln(m/3)$ at these prices, so
    $V(40) - V(30) = 3\ln(40/30) = 0.863$. The prediction overshoots because
    $\lambda^* = 3/m$ falls as $m$ rises: $V$ is concave, and a derivative
    taken at $m = 30$ overstates the average slope up to 40.
\end{enumerate}
\end{solution}

\begin{solution}{exo:b1:constrained-optimisation:corner}
The conditions $1 = \lambda$ and $1/y = \lambda$ give $y = 1$ and
$x = 0.5 - 1 = -0.5 < 0$: inadmissible. On the constraint,
$h(x) = x + \ln(0.5 - x)$ has $h'(x) = 1 - 1/(0.5 - x) < 0$ for
$0 \le x < 0.5$, so $h$ is largest at $x = 0$: the solution is the corner
$(0, 0.5)$. There the marginal value of $y$ is $1/y = 2$ and that of $x$ is
1: the ratio is 2, above the price ratio of 1, and the agent would like a
negative $x$ if it could have one.
\end{solution}

\begin{solution}{exo:b1:constrained-optimisation:student-kinks}
\begin{enumerate}
  \item The grant starts falling at $15h = 100$, $h = 6.67$, where income is
    $100 + 200 = \num{300}$~dalens. It is gone when $0.5(15h - 100) = 200$, that is
    $15h = 500$, $h = 33.33$, where income is \dal{500}.
  \item \dal{15}, then $15 \times (1 - 0.5) = \num{7.5}$~dalens, then \dal{15}.
  \item At 20 hours: earnings \dal{300}, grant $200 - 0.5 \times 200 = \num{100}$~dalens,
    income \dal{400}. At 40 hours: earnings \dal{600}, no grant, income \dal{600}.
\end{enumerate}
\end{solution}

\begin{solution}{exo:b1:constrained-optimisation:withdrawal}
With the rate 0.75 the grant is gone at $15h - 100 = 200/0.75 = 266.67$, that
is $h = 24.44$, and the middle segment pays $15 \times 0.25 = \num{3.75}$~dalens an
hour. Every point of the new constraint is on or below the old one, and the
first segment and the points beyond 33.33 hours are unchanged. The student at
5 hours chose that point when better points were available; it is still
available and the others are worse, so the student stays at 5 hours. The same
argument keeps the student at 40 hours (income there is still \dal{600}). The
student at 20 hours sees the wage on the margin halve (substitution: fewer
hours) and is poorer at 20 hours, \dal{350} instead of \dal{400} (income
effect: more hours); the direction is not settled by the model alone, but a
fall in hours down to the kink at 6.67 hours is the case the bunching
evidence points to. Only that student can end up at a kink.
\end{solution}

\begin{solution}{exo:b1:constrained-optimisation:equivalence}
\begin{enumerate}
  \item $\nabla g \neq 0$, so ``parallel'' means $\nabla f = \lambda \nabla g$ for
    some number $\lambda$: that is $f_x = \lambda g_x$, $f_y = \lambda g_y$, and
    with $g = b$ the \omterm{thm:b1:constrained-optimisation:lagrange}{Lagrange conditions} hold.
  \item The tangent lines at the point are perpendicular to $\nabla f$ and to
    $\nabla g$; they coincide if and only if the gradients are parallel, if and
    only if $f_x/f_y = g_x/g_y$, which is
    \cref{prop:b1:constrained-optimisation:equivalence}.
  \item Minimise $x^2 + y^2$ subject to $x + y = 2$: at $(1, 1)$,
    $\nabla f = (2, 2) = 2\nabla g$. The circles $x^2 + y^2 = c$ touch the line
    from the side of the origin; $(1, 1)$ is the point of the line closest to
    the origin, a \omterm{def:b1:constrained-optimisation:constrained-max}{constrained minimum}, and the \omterm{thm:b1:constrained-optimisation:lagrange}{Lagrange conditions} do not tell
    it from a maximum.
\end{enumerate}
\end{solution}

\begin{solution}{exo:b1:constrained-optimisation:dual}
\begin{enumerate}
  \item $\mathcal L = x + 2y - \mu(\sqrt{xy} - u_0)$. The conditions
    $1 = \tfrac{\mu}{2}\sqrt{y/x}$ and $2 = \tfrac{\mu}{2}\sqrt{x/y}$ give
    $x = 2y$; then $\sqrt{2}\,y = u_0$, so $y = u_0/\sqrt 2$,
    $x = \sqrt 2\, u_0$, $C(u_0) = 2\sqrt 2\, u_0$ and $\mu^* = 2\sqrt 2 = 2.828$.
  \item $C'(u_0) = 2\sqrt2 = \mu^*$, in dalens per unit of $f$: the cost of
    reaching one level higher.
  \item With $u_0 = \sqrt{18}$: $x = \sqrt{36} = 6$, $y = 3$, $C = 12$. The
    cheapest way to reach the level $\sqrt{18}$ costs exactly the budget whose
    best use reaches $\sqrt{18}$, at the same point. And
    $\mu^* = 2.828 = 1/0.354 = 1/\lambda^*$: one more unit of $f$ costs $1/\lambda^*$
    dalens when one more dalen buys $\lambda^*$ units.
\end{enumerate}
\end{solution}

\begin{solution}{exo:b1:constrained-optimisation:falsify}
\begin{enumerate}
  \item Nothing forces it to look different: the cloud compares countries, and
    countries differ in more than the wage, so even with every worker's supply
    rising in the wage the cloud could fall.
  \item For example: tax rates and the transfers withdrawn as earnings rise
    are higher in rich countries; their jobs are more often in offices with
    set hours, and paid holidays are longer.
  \item Hours of the same workers before and after a change in their wage
    that has nothing to do with their own choices (a change in a tax rate
    for one group only, say). If those workers' hours rose with the after-tax
    wage, the reading would be contradicted for them.
\end{enumerate}
\end{solution}

\begin{solution}{pb:b1:constrained-optimisation:farm}
\textbf{1.} $\Pi(x) = 2400\sqrt x + 1200\sqrt{20 - x}$ for $0 \le x \le 20$.

\textbf{2.} $\Pi'(x) = 1200/\sqrt x - 600/\sqrt{20 - x} = 0$ gives
$\sqrt{20 - x} = \sqrt x/2$, so $20 - x = x/4$, $x = 16$ and $y = 4$.

\textbf{3.} $\Pi''(x) = -600x^{-3/2} - 300(20 - x)^{-3/2} < 0$: $\Pi$ is
concave, so the critical point is the maximum.

\textbf{4.} $2400 \times 4 + 1200 \times 2 = \num{12000}$~dalens.

\textbf{5.} All potatoes: $2400\sqrt{20} = 2400 \times 4.4721 = \num{10733}$~dalens.
All barley: $1200 \times 4.4721 = \num{5367}$~dalens. Both are below \dal{12000}.

\textbf{6.} Potatoes: $1200/\sqrt{16} = \num{300}$~dalens a hectare; barley:
$600/\sqrt 4 = \num{300}$~dalens. If one were higher, moving a hectare to that crop
would raise profit, so the optimum cannot have them differ.

\textbf{7.} $\mathcal L = 2400\sqrt x + 1200\sqrt y - \lambda(x + y - 20)$.

\textbf{8.} $1200/\sqrt x = \lambda$, $600/\sqrt y = \lambda$, $x + y = 20$.

\textbf{9.} Dividing, $2\sqrt{y/x} = 1$, so $x/y = 4$.

\textbf{10.} $x = 4y$ and $5y = 20$: $y = 4$, $x = 16$, $\lambda = 1200/4 = 300$.

\textbf{11.} The level curve $2400\sqrt x + 1200\sqrt y = 12000$ is
$y = (10 - 2\sqrt x)^2$, convex towards the origin, and touches the line
$x + y = 20$ at $(16, 4)$, where both slopes are $-1$ ($-\Pi_x/\Pi_y = -300/300$).

\textbf{12.} Both find the critical point of $\Pi$ along the constraint:
substitution by solving the constraint, Lagrange by the \omterm{def:b1:constrained-optimisation:tangency}{tangency condition},
which \cref{prop:b1:constrained-optimisation:equivalence} shows is the same
condition. $\Pi$ is concave, so that point is the maximum.

\textbf{13.} The \omterm{def:b1:constrained-optimisation:tangency}{tangency condition} still gives $x = 4y$, so $x = 4L/5$ and
$y = L/5$. Then $V(L) = 2400 \times 2\sqrt{L/5} + 1200\sqrt{L/5} = 6000\sqrt{L/5}$.

\textbf{14.} $V'(L) = 6000 \times \tfrac{1}{2\sqrt{L/5}} \times \tfrac15 =
600/\sqrt{L/5}$; at $L = 20$, $600/2 = 300 = \lambda$
(\cref{prop:b1:constrained-optimisation:shadow}).

\textbf{15.} Dalens of profit per hectare for the season: what one more
hectare is worth to this farm, the most it would pay to rent one.

\textbf{16.} At \dal{250}, yes: the hectare is worth about \dal{300}. At
\dal{350}, no.

\textbf{17.} $V(21) - V(20) = 6000(\sqrt{4.2} - 2) = 6000 \times 0.04939 =
\dal{296.34}$, a little below 300 because $V$ is concave: the shadow rent
falls as the field grows.

\textbf{18.} It falls: $V'(L) = 600/\sqrt{L/5}$ decreases in $L$. Each extra
hectare goes to poorer soil.

\textbf{19.} $16 \times 500 + 4 \times 100 = 8400$ cubic metres. With
$W = 9000$ the water constraint is slack and its \omterm{def:b1:constrained-optimisation:shadow-price}{shadow price} is zero.

\textbf{20.} The plan of Part II needs 8400 cubic metres, more than 6000, so
water binds. If land were slack, the water constraint alone would give
$1200/\sqrt x = 500\mu$ and $600/\sqrt y = 100\mu$, so $y = 6.25x$, and
$1125x = 6000$ gives $y = 33.3 > 20$: land would be exceeded, so it binds
too. From $x + y = 20$ and $500x + 100y = 6000$: $400x + 2000 = 6000$,
$x = 10$, $y = 10$.

\textbf{21.} $2400\sqrt{10} + 1200\sqrt{10} = 3600 \times 3.1623 = \num{11384}$~dalens;
the water allowance costs $12000 - 11384 = \num{616}$~dalens.

\textbf{22.} $\mathcal L = 2400\sqrt x + 1200\sqrt y - \lambda(x + y - 20) -
\mu(500x + 100y - 6000)$, with conditions $1200/\sqrt x = \lambda + 500\mu$,
$600/\sqrt y = \lambda + 100\mu$, and the two constraints.

\textbf{23.} At $x = y = 10$: $379.47 = \lambda + 500\mu$ and
$189.74 = \lambda + 100\mu$. Subtracting, $400\mu = 189.73$, so
$\mu = 0.4743$; then $\lambda = 189.74 - 47.43 = 142.31$.

\textbf{24.} With plentiful water, one more hectare is worth \dal{300} to the
farm; with an allowance of 6000 cubic metres, about \dal{142}, because a new
hectare must be planted mostly in barley to stay within the water. One more
cubic metre of water is worth about \dal{0.47}.
\end{solution}
